%% kapa: exam=1 round=26 year=2015 period=1 form=A subject=economics subject_ko=경제원론 seq=10 number=50
%% source: Q-Net 감정평가사 기출문제 (cst00309) artlSeq=5138001 file=extracted/1008/5138001/제26회_감정평가사 1차_기출/제26회_감정평가사 1차_기출/1교시_감평_A형.pdf
%% conversion: kordoc 4.24.51 -> tools/gen1.py (auto); review=proofread-sonnet; flags=none
\begin{question}{50}
甲기업의 생산함수는 $f(K,L)=K^{1/2}L^{1/4}$이고, 산출물의 가격은 4, $K$의 가격은 2, $L$의 가격은 1이다. 이윤을 극대화하는 甲기업의 $K$와 $L$은 각각 얼마인가? (단, $K$와 $L$은 각각 자본, 노동 투입량을 나타내고, 생산물시장과 생산요소시장은 완전경쟁시장이다.)
\begin{choices}
  \item $K=1, L=1$
  \item $K=1, L=2$
  \item $K=2, L=2$
  \item $K=2, L=4$
  \item $K=4, L=2$
\end{choices}
\end{question}
\answer{1}
